\documentclass[10pt,a4paper]{amsart}
\usepackage[margin=19mm]{geometry}
\usepackage{amsmath,amssymb,mathtools}
\usepackage[scr=rsfs,cal=euler]{mathalfa}
\usepackage{xfrac}
\usepackage[unicode,hidelinks,bookmarks=false]{hyperref}
\newcommand{\ie}{i.e.\ }
\newcommand{\cf}{cf.\ }
\newcommand{\resp}{respectively\ }
\newcommand{\Subsection}{\subsection}
\providecommand{\nobreakdash}{\nobreak\hbox{-}}
\setlength{\emergencystretch}{2em}
\multlinegap0pt
\usepackage{amssymb}
\usepackage{mathtools}
\usepackage[shortlabels]{enumitem}
\usepackage{xcolor}
\newtheorem{theorem}{Theorem}
\newtheorem{proposition}{Proposition}
\usepackage{cleveref}
\crefname{theorem}{Theorem}{Theorems}
\crefname{proposition}{Proposition}{Propositions}
\newcommand{\F}{\mathbb F}
\newcommand{\SL}{\operatorname{SL}}
\newcommand{\cE}{\mathcal E}
\newcommand{\cN}{\mathcal N}
\newcommand{\cS}{\mathcal S}
\title{Lower bounds for some value sets over finite fields:\\ incidence geometry and Bourgain's group expansion theorem}
\author{Xiyu Hu}
\date{Source: arXiv:2609.05652v1 (CC BY 4.0)}
\begin{document}
\maketitle

\noindent\emph{Source and licence.} Lower bounds for some value sets over finite fields:\\ incidence geometry and Bourgain's group expansion theorem. Version arXiv:2609.05652v1 (CC BY 4.0), distributed by arXiv under the Creative Commons Attribution 4.0 International licence, http://creativecommons.org/licenses/by/4.0/, as recorded in the arXiv licence metadata for this exact version and shown on the abstract page https://arxiv.org/abs/2609.05652v1. Full licence notice: \texttt{licenses/arxiv-2609.05652-CC-BY-4.0.txt}.\par\smallskip

\noindent\textit{Editorial scope.} Counted unit: the article's theorem thm:abstract-intro (source0.tex lines 276--281). The article's complete original proof of it is delivered with it (source0.tex lines 947--1005, one \texttt{proof} environment of 59 lines, which the article places away from the statement). No mathematics is written, repaired or re-derived: the statement and the proof are the article's own lines, in the article's own notation, and no retained line is substituted or re-flowed.  Retained uncounted as the source-local context the proof needs: lines 258--264; lines 755--766; lines 843--845; lines 892--901.  Nothing was omitted from any retained span, so the package carries no omission record.  No retained line contains a citation, so the wrapper carries no bibliography and no bibliography entry.  No retained line contains a figure environment, an image-inclusion command or drawing code, so no image is delivered and the package declares no asset dependency.  Editorial formatting only, in addition: an AMS \texttt{amsart} preamble that restates the article's own declarations and macros used by the retained lines, namely the environments \texttt{theorem}, \texttt{proposition} on separate counters, together with the standard packages those lines need.  The retained environments print in this document as Theorem 1, Proposition 1 in order of appearance, whereas the article prints the same items with the numbering of its own full document.  The complete native source of the article as distributed in the arXiv e-print for this exact version is bundled unchanged as \texttt{sources/209454/source0.tex}.
\begin{equation}\label{eq:phi-intro}
  \varphi(t)=\frac{\alpha t+\beta}{\gamma t+\delta},
  \qquad
  \alpha,\beta,\gamma,\delta\in\F_p^\times,
  \qquad
  \Delta:=\alpha\delta-\beta\gamma\ne0,
\end{equation}

\begin{theorem}[Bourgain]\label{thm:bourgain}
For every $\varepsilon>0$ and $r>1$, there is a number $\sigma=\sigma(\varepsilon,r)>0$ such that the following holds for all sufficiently large primes $p$.  Let $A\subseteq\F_p$ and $S\subseteq\SL_2(\F_p)$ satisfy
\begin{align}
  1\ll |A|&<p^{1-\varepsilon},\label{eq:bourgain-size-A}\\
  \log|A|&<r\log|S|,\label{eq:bourgain-size-S}\\
  |S\cap gK|&<|S|^{1-\varepsilon}\label{eq:bourgain-nonconc}
\end{align}
for every proper subgroup $K<\SL_2(\F_p)$ and every $g\in\SL_2(\F_p)$.  Then
\begin{equation}\label{eq:bourgain-conclusion}
  I(A,S)\ll |A|^{1-\sigma}|S|.
\end{equation}
\end{theorem}

\begin{equation}\label{eq:S-size-exact}
  |\cS_V|=m(m-1).
\end{equation}

\begin{proposition}\label{prop:subgroup-escape}
There is an absolute constant $C$ such that, for every prime $p\ge5$, every $V\subseteq\F_p^\times$ with $m=|V|$, every proper subgroup $K<\SL_2(\F_p)$, and every $g\in\SL_2(\F_p)$,
\begin{equation}\label{eq:subgroup-escape}
  |\cS_V\cap gK|\le Cm+C.
\end{equation}
In particular, for all sufficiently large $m$,
\begin{equation}\label{eq:subgroup-escape-power}
  |\cS_V\cap gK|<|\cS_V|^{3/4}.
\end{equation}
\end{proposition}

\begin{theorem}\label{thm:abstract-intro}
There are absolute constants $c,\eta>0$ such that, for every prime $p$, every quadruple satisfying \eqref{eq:phi-intro}, and every $V\subseteq\F_p^\times$,
\begin{equation}\label{eq:abstract-intro}
  |V|\ge c\min\{\cN_T(V),p\}^{1/2+\eta}.
\end{equation}
\end{theorem}

\begin{proof}[Proof of \cref{thm:abstract-intro}]
Write $m=|V|$.  By Cauchy--Schwarz,
\begin{equation}\label{eq:CS-main}
  L^2
  \le m\sum_{x\in V}R_{\cE}(x)^2
  \le m\sum_{a,b\in V}M(a,b),
\end{equation}
where
\begin{equation}\label{eq:M-ab}
  M(a,b)
  =\#\{x\in V:T_a(x)\in V,\ T_b(x)\in V\}.
\end{equation}
The diagonal terms contribute at most $m^2$.

Suppose $a\ne b$.  Put $u=T_a(x)$.  Since $T_a$ is a projective automorphism, the substitution is injective, and
\[
  T_b(x)=(T_b\circ T_a^{-1})(u)=H_{a,b}\cdot u.
\]
Consequently,
\begin{equation}\label{eq:M-incidence}
  \sum_{\substack{a,b\in V\\a\ne b}}M(a,b)
  \le I(V,\cS_V).
\end{equation}

Fix $\varepsilon=1/4$ and $r=2$ in \cref{thm:bourgain}, and let $\sigma>0$ be the resulting saving.  We may assume $0<\sigma\le1$.

If $m\ge p^{3/4}$, then $L\le p$ gives
\begin{equation}\label{eq:large-m-case}
  m\ge L^{3/4}.
\end{equation}
Suppose therefore that $m<p^{3/4}$.  For sufficiently large $m$, \eqref{eq:S-size-exact} gives
\[
  \log m<2\log|\cS_V|,
\]
and \cref{prop:subgroup-escape} gives
\[
  |\cS_V\cap gK|<|\cS_V|^{3/4}
\]
for every proper subgroup $K$ and every coset.  Thus all hypotheses of \cref{thm:bourgain} are satisfied, and
\begin{equation}\label{eq:I-bound}
  I(V,\cS_V)
  \ll m^{1-\sigma}|\cS_V|
  \ll m^{3-\sigma}.
\end{equation}
Combining \eqref{eq:CS-main}, \eqref{eq:M-incidence} and \eqref{eq:I-bound}, we obtain
\begin{equation}\label{eq:L-bound}
  L^2\ll m(m^2+m^{3-\sigma})\ll m^{4-\sigma}.
\end{equation}
Hence
\begin{equation}\label{eq:m-from-L}
  m\gg L^{2/(4-\sigma)}.
\end{equation}
Set
\begin{equation}\label{eq:eta-def}
  \eta=\frac{2}{4-\sigma}-\frac12
  =\frac{\sigma}{2(4-\sigma)}>0.
\end{equation}
The exponent in \eqref{eq:large-m-case} is stronger than $1/2+\eta$.  The finitely many cases in which $p$ or $m$ is below the thresholds in Bourgain's theorem are absorbed by decreasing the absolute implied constant.  Taking $L=\min\{\cN_T(V),p\}$ proves \eqref{eq:abstract-intro}.
\end{proof}

\end{document}
